\documentclass[11pt]{amsart}

\usepackage[T1]{fontenc}
\usepackage{lmodern}
\usepackage{microtype}
\usepackage{amsmath,amssymb,amsthm}
\usepackage[hidelinks]{hyperref}

\hypersetup{
  pdftitle={Small Counterexamples to a Trace Conjecture of Frankl and Wang}
}

\newtheorem{theorem}{Theorem}
\newtheorem{proposition}[theorem]{Proposition}

\newcommand{\tr}{\mathbin{|}}
\newcommand{\cF}{\mathcal F}
\newcommand{\cG}{\mathcal G}
\newcommand{\cH}{\mathcal H}
\newcommand{\cL}{\mathcal L}
\newcommand{\cM}{\mathcal M}
\newcommand{\cQ}{\mathcal Q}

\title[Small counterexamples to a trace conjecture]
{Small Counterexamples to a Trace Conjecture of Frankl and Wang}
\date{}
\subjclass[2020]{Primary 05D05; Secondary 05B07}
\keywords{trace of a set family, arrow relation, extremal set theory,
Kruskal--Katona theorem, Fano plane}

\begin{document}

\begin{abstract}
Frankl and Wang conjectured an exact formula for
$m(n,\ell+1,3\cdot2^{\ell-1}+1)$ for every $n>\ell>0$.
We prove
\[
  m(6,5,25)=40,\qquad m(7,5,25)\ge58,\qquad
  m(7,6,49)\ge80,
\]
where their formula gives $37$, $55$, and $73$, respectively.
Thus the conjecture fails already at $n=\ell+2$ for $\ell=4$ and $5$.
The first value is determined using the Kruskal--Katona theorem; the
other counterexamples are explicit down-sets.
\end{abstract}

\maketitle

\section{Introduction}

For $\cF\subseteq2^{[n]}$ and $Y\subseteq[n]$, let
\[
  \cF\tr_Y=\{F\cap Y:F\in\cF\}.
\]
The arrow relation $(n,m)\to(a,b)$ means that every
$\cF\subseteq2^{[n]}$ with $|\cF|\ge m$ has an $a$-set $Y$ satisfying
$|\cF\tr_Y|\ge b$. Write
\[
  m(n,a,b)=\min\{m:(n,m)\to(a,b)\}.
\]
Frankl and Wang conjectured \cite[Conjecture~1.3]{FW} that, for all
$n>\ell>0$,
\begin{equation}\label{eq:FW}
  m\bigl(n,\ell+1,3\cdot2^{\ell-1}+1\bigr)
  =1+\prod_{0\le i<\ell}
  \left\lfloor\frac{n+\ell+i}{\ell}\right\rfloor.
\end{equation}
The same statement was subsequently recorded as an open problem in
\cite[Conjecture~2]{LMR}.

\begin{theorem}\label{thm:main}
One has
\[
  m(6,5,25)=40,\qquad m(7,5,25)\ge58,
  \qquad m(7,6,49)\ge80.
\]
Consequently, \eqref{eq:FW} is false for
$(n,\ell)=(6,4),(7,4),(7,5)$.
\end{theorem}

For these parameters, the right-hand side of \eqref{eq:FW} is,
respectively, $37$, $55$, and $73$. When $n=\ell+1$, the proposed
threshold equals the target trace size, so the assertion is tautological.
Thus $(6,4)$ and $(7,5)$ are the first nontrivial cases for $\ell=4$
and $5$.

All families constructed below are down-sets. For a down-set $\cF$,
\begin{equation}\label{eq:downtrace}
  \cF\tr_Y=\cF\cap2^Y.
\end{equation}
We write $ijk$ for $\{i,j,k\}$ and
$\binom{X}{\le r}=\{S\subseteq X:|S|\le r\}$.

\section{The exact value
\texorpdfstring{$m(6,5,25)=40$}{m(6,5,25)=40}}

Let
\[
  D=\{136,145,156,234,236,245\},\qquad
  A=\{1235,1246,3456\},
\]
and define
\begin{equation}\label{eq:F6}
  \cF=\binom{[6]}{\le2}
  \cup\left(\binom{[6]}3\setminus D\right)\cup A.
\end{equation}
No member of $D$ is contained in a member of $A$, so $\cF$ is a
down-set. Moreover, $D$ is $3$-regular and $A$ is $2$-regular on
$[6]$, while
\[
  |\cF|=22+(20-6)+3=39.
\]
For $Y_x=[6]\setminus\{x\}$, equation \eqref{eq:downtrace} gives
\[
  |\cF\tr_{Y_x}|
  =\sum_{i=0}^{2}\binom5i+\bigl(\tbinom53-(6-3)\bigr)+(3-2)
  =16+7+1=24.
\]
Hence $m(6,5,25)\ge40$.

\begin{proposition}\label{prop:upper}
Every $40$-member family in $2^{[6]}$ has a $5$-element trace of size
at least $25$.
\end{proposition}

\begin{proof}
Suppose otherwise. By the standard squashing lemma
\cite[Lemma~1.1]{FW}, the family may be replaced by a $40$-member
down-set with no larger traces. Denote this down-set by $\cF$ and put
$\cG=2^{[6]}\setminus\cF$. Then $\cG$ is an up-set of size $24$.
For $Y_x=[6]\setminus\{x\}$, the assumption and
\eqref{eq:downtrace} imply
\begin{equation}\label{eq:eight}
  |\cG\cap2^{Y_x}|=32-|\cF\tr_{Y_x}|\ge8
  \qquad(x\in[6]).
\end{equation}

The family $\cG$ contains no set of size at most one, since the upper
cone of such a set has at least $32$ members. It also contains no pair.
Indeed, suppose $\{u,v\}\in\cG$. Its upper cone has $16$ members, so
exactly eight members of $\cG$ lie outside that cone. Equation
\eqref{eq:eight} for $x=u$ forces all eight to avoid $u$, and the same
equation for $x=v$ forces all eight to avoid $v$. If $E$ is one of
them, then $E\cup\{u\}\in\cG$ because $\cG$ is an up-set. This set
still avoids $v$, so it lies outside the cone, but it contains $u$, a
contradiction.

Thus every member of $\cG$ has size at least three. For each $x$,
equation \eqref{eq:eight} gives some member of $\cG$ contained in
$Y_x$; the up-set property then gives $Y_x\in\cG$. Hence all six
$5$-sets and $[6]$ belong to $\cG$. Set
\[
  t=\left|\cG\cap\binom{[6]}3\right|,
  \qquad
  q=\left|\cG\cap\binom{[6]}4\right|.
\]
Then $t+q=17$, and double counting pairs $(x,G)$ with
$G\in\cG\cap2^{Y_x}$ yields
\[
  48\le\sum_{x=1}^{6}|\cG\cap2^{Y_x}|=3t+2q+6.
\]
Consequently, $t\ge8$ and $q\le9$.

Choose eight triples from $\cG$ and call the resulting family $\cH$.
The $4$-element upper shadow of $\cH$ is contained in
$\cG\cap\binom{[6]}4$. Under complementation in $[6]$, this upper
shadow becomes the ordinary $2$-element lower shadow of the eight
triples $\{[6]\setminus H:H\in\cH\}$. Since
\[
  8=\binom43+\binom32+\binom11,
\]
the Kruskal--Katona theorem \cite{Katona,Kruskal} gives
\[
  q\ge\binom42+\binom31+\binom10=10,
\]
contrary to $q\le9$.
\end{proof}

The construction \eqref{eq:F6} and Proposition~\ref{prop:upper} prove
$m(6,5,25)=40$.

\section{Two counterexamples on seven points}

Let
\[
  \cL=\{123,145,167,246,257,347,356\}
\]
be the lines of a Fano plane on $[7]$, and set
\[
  \cF_4=\binom{[7]}{\le3}\setminus\cL.
\]
Then $|\cF_4|=64-7=57$. Every $5$-set is the complement of a pair.
Exactly five Fano lines meet that pair, so exactly two Fano lines lie
inside the $5$-set. Therefore every $5$-trace has size
\[
  \sum_{i=0}^{2}\binom5i+\binom53-2=24,
\]
and hence
\begin{equation}\label{eq:m7525}
  m(7,5,25)\ge58.
\end{equation}

For the final construction, let
\[
  C=\{1,4\},\qquad
  \cM=\{123,456,147\},
\]
and define
\[
  \cQ=\left\{Q\in\binom{[7]}4:
  Q\cap C\ne\varnothing\ \text{and}\ M\nsubseteq Q
  \text{ for every }M\in\cM\right\}.
\]
There are $35-\binom54=30$ four-sets meeting $C$. Each member of
$\cM$ has four four-set supersets, all meeting $C$, and these three
upper cones are disjoint because the union of any two members of
$\cM$ has size at least five. Thus $|\cQ|=30-12=18$.

Now set
\begin{equation}\label{eq:F7}
  \cF_5=\binom{[7]}{\le2}
  \cup\left(\binom{[7]}3\setminus\cM\right)\cup\cQ.
\end{equation}
By the definition of $\cQ$, this is a down-set, and
\[
  |\cF_5|=29+(35-3)+18=79.
\]
Let $Y_x=[7]\setminus\{x\}$. If $x\in C$, exactly one member of
$\cM$ lies in $Y_x$. Moreover, $Y_x$ contains ten four-sets meeting
$C$, exactly three of which contain that member of $\cM$; hence
$|\cQ\cap2^{Y_x}|=7$. Therefore
\[
  |\cF_5\tr_{Y_x}|=22+(20-1)+7=48.
\]
If $x\notin C$, exactly two members of $\cM$ lie in $Y_x$. Among the
fifteen four-sets in $Y_x$, exactly one is disjoint from $C$. Each of
the two contained members of $\cM$ has three four-set supersets in
$Y_x$, and no four-set contains both. Hence
$|\cQ\cap2^{Y_x}|=14-6=8$, and
\[
  |\cF_5\tr_{Y_x}|=22+(20-2)+8=48.
\]
Thus
\begin{equation}\label{eq:m7649}
  m(7,6,49)\ge80.
\end{equation}
Equations \eqref{eq:m7525} and \eqref{eq:m7649} complete the proof of
Theorem~\ref{thm:main}.

Theorem~\ref{thm:main} disproves \eqref{eq:FW} as stated. It does not
address whether the same formula might hold for all sufficiently large
$n$ when $\ell$ is fixed.

\begin{thebibliography}{9}

\bibitem{FW}
P.~Frankl and J.~Wang,
\emph{Four-vertex traces of finite sets},
Graphs Combin. \textbf{40} (2024), Art.~8,
\href{https://doi.org/10.1007/s00373-023-02738-5}
{doi:10.1007/s00373-023-02738-5}.

\bibitem{Katona}
G.~O.~H.~Katona,
\emph{A theorem of finite sets},
in Theory of Graphs (Proc. Colloq., Tihany, 1966),
Academic Press, New York, 1968, pp.~187--207.

\bibitem{Kruskal}
J.~B.~Kruskal,
\emph{The number of simplices in a complex},
in Mathematical Optimization Techniques,
University of California Press, Berkeley, 1963, pp.~251--278.

\bibitem{LMR}
M.~Li, J.~Ma, and M.~Rong,
\emph{Recent advances in arrow relations and traces of sets},
in Sum(m)it280, Bolyai Soc. Math. Stud. \textbf{32},
Springer, Cham, 2026, pp.~327--354,
\href{https://doi.org/10.1007/978-3-032-18810-6_14}
{doi:10.1007/978-3-032-18810-6\_14}.

\end{thebibliography}

\end{document}